% !TEX program = xelatex
\documentclass[aspectratio=169,9pt]{beamer}

% EPIC Lab Beamer theme. Keep this .tex file next to beamerthemeEPIC.sty.
% Do not change aspectratio: the theme's frametitle rule and section dividers
% use hard-coded tikz coordinates tuned for 16:9.
\usetheme{EPIC}

\usepackage{amsmath,amssymb,mathtools}
\usepackage{centernot}
\usepackage{booktabs}
\usepackage{array}
\usepackage{tikz}
\usetikzlibrary{calc,arrows.meta,positioning,fit}

% Semantic highlighting. \orange = the key technical term of a bullet,
% \blue = a secondary or linking concept, \muted = de-emphasised aside.
\newcommand{\orange}[1]{\textcolor{AmazonOrange}{\textbf{#1}}}
\newcommand{\blue}[1]{\textcolor{densityblue}{\textbf{#1}}}
\newcommand{\green}[1]{\textcolor{deepgreen}{\textbf{#1}}}
\newcommand{\muted}[1]{\textcolor{EPICMuted}{#1}}
\definecolor{densityblue}{RGB}{0,75,155}
\definecolor{deepgreen}{RGB}{20,130,95}

% Numbered statements. Blocks carrying a formal statement are labelled by hand as
% "Definition <section>.<n>", "Proposition <section>.<n>" or "Theorem <section>.<n>",
% numbered consecutively within each \section. Numbers are literal rather than
% counter-driven: a LaTeX counter stepped inside a frame is re-stepped on every
% \pause overlay, which would corrupt the numbering. Never leave a block title empty.
% Shorthands for the information-theoretic quantities used throughout.
\newcommand{\KL}[2]{D_{\mathrm{kl}}\!\left(#1 \,\|\, #2\right)}
\newcommand{\E}{\mathbb{E}}
\newcommand{\Prob}{\mathbb{P}}

% Node styles for the source/channel coding pipeline diagrams.
\tikzset{
  pipebox/.style={draw=OxfordBlue,fill=OxfordBlue!6,rounded corners=2pt,
                  inner xsep=5pt,inner ysep=4pt,font=\footnotesize,minimum height=0.62cm},
  pipechan/.style={draw=AmazonOrange,fill=AmazonOrange!15,rounded corners=2pt,
                   inner xsep=5pt,inner ysep=4pt,font=\footnotesize\bfseries,minimum height=0.62cm},
  pipearr/.style={-{Latex[length=1.7mm]},OxfordBlue,thick},
}

% Citations. Uncomment this block together with the References frame below
% once reference.bib has entries.
% \usepackage{natbib}
% \setbeamertemplate{bibliography item}[text]

% The optional short title is required: the footer prints
% \insertshorttitle | frame/total, and degrades silently without it.
\title[Information Theory Reading Group]{Information Theory Reading Group}
\subtitle{Duchi, \textit{Lecture Notes on Statistics and Information Theory}, Chapters 1 and 2.1}
\author{Joshua Adrian Cahyono}
\institute{EPIC LAB Nanyang Technological University}
\date{14 September 2026}

\begin{document}

% ------------------------------------------------------------------
\begin{frame}[plain]
  \titlepage
\end{frame}

% ------------------------------------------------------------------
% Talk structure. A single column of four, set large and generously spaced so the
% roadmap fills the slide rather than hugging the frame title.
\begin{frame}{Talk structure}
\vfill
{\LARGE
\begin{enumerate}
    \setlength{\itemsep}{0.5cm}
    \item What is information theory?
    \item Entropy: measuring uncertainty.
    \item KL divergence: the cost of a wrong model.
    \item Mutual information: what data tells us.
\end{enumerate}
}
\vfill
\end{frame}

\section{What is Information Theory?}
% \section{} automatically generates a full-bleed divider frame via the
% theme's \AtBeginSection hook. Never hand-roll a divider slide.

% ------------------------------------------------------------------
\begin{frame}{Scope of the subject}
\epicquote{Information and uncertainty are two sides of the same coin.}{C. E. Shannon}
\vspace{0.2cm}
\pause
\begin{block}{Working definition}
Information theory studies how to \orange{measure uncertainty}, and how much of it we \orange{remove by observing data}.
\end{block}
\pause
\vspace{0.1cm}
It gives us three things, which we build up over the coming sessions:
\begin{itemize}
    \item a way to put uncertainty on a scale, in \orange{bits};
    \item a way to measure how far apart two distributions are, called a \orange{divergence};
    \item results telling us \orange{the best that any method could possibly do}, however clever it is.
\end{itemize}
\end{frame}

% ------------------------------------------------------------------
\begin{frame}{Origins: Shannon (1948)}
Three choices Shannon made, which still shape the subject today.
\begin{itemize}
    \item \orange{Meaning is set aside.} What matters about a message is how unexpected it was, not what it says:
    \textit{``the semantic aspects of communication are irrelevant to the engineering problem.''}
    \pause
    \item \orange{The source is treated as random.} Once a message is a random variable, we can use probability
    to study it. This is the step that ties information theory to statistics.
    \pause
    \item \orange{The limits can actually be reached.} Shannon showed not only that a certain rate cannot be
    beaten, but that we can get as close to it as we like.
\end{itemize}
\vspace{0.3cm}
\pause
\muted{The last point is what makes the theory useful rather than discouraging: its bounds are targets we can
aim at, not just bad news.}
\end{frame}

% ------------------------------------------------------------------
\begin{frame}{Information and uncertainty}
How much a message tells us depends on \blue{how likely the event was}, and on nothing else:
\begin{center}
\large information $\;=\;$ reduction in uncertainty $\;=\;$ surprise.
\end{center}
\vspace{0.1cm}
\pause
\begin{center}
\footnotesize
\begin{tabular}{lcc}
\toprule
\textbf{Message} & $\Prob(\text{event})$ & \textbf{How much it tells us} \\
\midrule
``The sun rose this morning.'' & $\approx 1$ & $\approx 0$ bits \\
``It snowed in Singapore this morning.'' & small & a great deal \\
\bottomrule
\end{tabular}
\end{center}
\vspace{0.3cm}
\pause
\begin{block}{Remark}
We always measure against \textit{what we expected}, so the same message tells different people different
amounts. This is why the uncertainty in a model's prediction describes \blue{the model}, not the world.
We return to this in Section 4.
\end{block}
\end{frame}

% ------------------------------------------------------------------
\begin{frame}{Units: the bit}
\begin{block}{Definition 1.1 (informal)}
One \orange{bit} is the uncertainty in a single fair coin flip; equally, it is what we learn from the answer to
one well chosen yes/no question.
\end{block}
\pause
\vspace{0.15cm}
Finding one item among eight equally likely ones takes three yes/no questions, and
\[
    \log_2 8 = 3 .
\]
\pause
\vspace{-0.15cm}
\begin{itemize}
    \item The logarithm appears because the number of questions grows with the \textit{log} of the number of possibilities.
    \pause
    \item We make this exact twice in Part 2: once for compression, and once for the twenty questions game.
\end{itemize}
\end{frame}

% ------------------------------------------------------------------
\begin{frame}{The two problems Shannon posed \muted{(Duchi, \S1.1)}}
\begin{block}{Problem 1. Compression: how few bits can represent the signal?}
\centering
\begin{tikzpicture}[node distance=0.5cm]
    \node[pipebox] (s) {Source};
    \node[pipebox,right=of s] (c) {Compressor};
    \node[pipebox,right=of c] (d) {Decompressor};
    \node[pipebox,right=of d] (r) {Receiver};
    \draw[pipearr] (s) -- (c); \draw[pipearr] (c) -- (d); \draw[pipearr] (d) -- (r);
\end{tikzpicture}
\end{block}
\pause
\begin{block}{Problem 2. Transmission: how much can we send over a noisy link?}
\centering
\begin{tikzpicture}[node distance=0.5cm]
    \node[pipebox] (s) {Source};
    \node[pipebox,right=of s] (c) {Compressor};
    \node[pipechan,right=of c] (ch) {Channel};
    \node[pipebox,right=of ch] (d) {Decompressor};
    \node[pipebox,right=of d] (r) {Receiver};
    \draw[pipearr] (s) -- (c); \draw[pipearr] (c) -- (ch);
    \draw[pipearr] (ch) -- (d); \draw[pipearr] (d) -- (r);
\end{tikzpicture}
\end{block}
\pause
\vspace{0.1cm}
\muted{In the notes these are called \textit{source coding} and \textit{channel coding}.}
The first is answered by \orange{entropy}, the second by \orange{mutual information}.
The rest of Chapter 2 builds the tools for both.
\end{frame}

% ------------------------------------------------------------------
\begin{frame}{The same two problems in statistics \muted{(Duchi, \S1.2)}}
\begin{center}
\begin{tikzpicture}[node distance=0.9cm]
    \node[pipebox] (s) {Source $P$};
    \node[pipebox,right=of s] (c) {Compressor};
    \node[pipechan,right=of c] (ch) {Channel $P(Y \mid f(X))$};
    \node[pipebox,right=of ch] (d) {Decompressor};
    \draw[pipearr] (s) -- node[above,font=\tiny]{$X_{1:n}$} (c);
    \draw[pipearr] (c) -- node[above,font=\tiny]{$f(X_{1:n})$} (ch);
    \draw[pipearr] (ch) -- node[above,font=\tiny]{$Y_{1:n}$} (d);
\end{tikzpicture}
\end{center}
\vspace{0.1cm}
\begin{itemize}
    \item \orange{Compression becomes modelling.} Given samples $X_1,\dots,X_n$ from $P$, we build a model
    $\widehat P$ that describes them compactly: \blue{density estimation and generative modelling}.
    \pause
    \item \orange{Transmission becomes inference.} Nature picks $\theta$, we see a noisy $Y$, and we have to
    recover $\theta$: \blue{estimation and inference}.
    \pause
    \item \orange{The key difference.} In communication we get to design the compressor. In statistics the way
    the data were collected hands us one, and we are stuck with it.
    \pause
    \item \muted{Duchi lists the cases where we get some of that freedom back: experimental design, active
    learning, bandits, reinforcement learning.}
\end{itemize}
\vspace{0.2cm}
\pause
\takeaway{Statistics and machine learning are Shannon's two problems, asked of someone who does not control the encoder.}
\end{frame}

% ------------------------------------------------------------------
\begin{frame}{Limits that hold for every method}
\begin{block}{The general form}
Information theory proves statements of the form: \textit{no method of this kind can do better than this.}
\end{block}
\pause
\vspace{0.2cm}
\begin{itemize}
    \item Results like these separate a weakness in \textit{one particular model} from a limit built into
    \textit{the problem itself}.
    \pause
    \item That makes them practical. Such a bound tells us whether it is worth tuning further, or whether we
    need better data or a better representation instead.
    \pause
    \item Chapter 2 proves two of them, due to Le Cam and to Fano. Both are the business of Part 2.
\end{itemize}
\vspace{0.3cm}
\pause
\muted{These are the results the notes call \textit{converse results}. They are the main reason the subject is
worth knowing outside communications, and the thread we follow throughout.}
\end{frame}

% ------------------------------------------------------------------
\begin{frame}{What the theory does not cover}
\begin{itemize}
    \item \orange{Meaning or usefulness.} A completely random string carries the most information in Shannon's
    sense, and is worth nothing. Whether information is \textit{useful} is a question for decision theory.
    \pause
    \item \orange{Entropy in physics.} The connection to thermodynamics is real, but we do not pursue it here.
    \pause
    \item \orange{Signals on wires.} The theory needs only uncertainty and observation, which is why the same few
    quantities turn up in statistics, neuroscience and machine learning.
\end{itemize}
\vspace{0.25cm}
\pause
\begin{alertblock}{A caution about words}
Here, ``high information'' means ``unlikely, given what we expected''. It does not mean ``important''.
\end{alertblock}
\end{frame}

% ------------------------------------------------------------------
\begin{frame}{The quantities we will meet}
\begin{center}
\begin{tabular}{clp{5.8cm}c}
\toprule
\textbf{Symbol} & \textbf{Name} & \textbf{What it asks} & \textbf{Section} \\
\midrule
$H(X)$            & entropy            & how uncertain are we about $X$?                & 2 \\[0.12cm]
$\KL{P}{Q}$       & KL divergence      & what does it cost to use $Q$ in place of $P$?  & 3 \\[0.12cm]
$I(X;Y)$          & mutual information & how much does $Y$ tell us about $X$?           & 4 \\[0.12cm]
$D_f(P\|Q)$       & $f$-divergence     & the general family containing the others       & 5 \\
\bottomrule
\end{tabular}
\end{center}
\vspace{0.45cm}
\pause
Several familiar things in machine learning are these quantities under other names:
\begin{itemize}
    \item cross-entropy loss is a KL divergence, plus a constant that does not depend on the model;
    \item perplexity is the entropy, exponentiated;
\end{itemize}
\end{frame}

% ------------------------------------------------------------------
\begin{frame}{Prerequisites}
\begin{columns}[T,totalwidth=\textwidth]
\begin{column}{0.55\textwidth}
\begin{itemize}
    \item \orange{Expectation}, a weighted average,
    \[
        \E_P[f(X)] = \sum_x p(x) f(x).
    \]
    \pause
    \item \orange{Convexity}: the straight line joining two points on the curve stays above it,
    \[
        f(\lambda x + (1-\lambda)y) \le \lambda f(x) + (1-\lambda) f(y).
    \]
    \pause
    \item \orange{Jensen's inequality}, which follows at once,
    \[
        f(\E[X]) \le \E[f(X)].
    \]
\end{itemize}
\end{column}
\begin{column}{0.4\textwidth}
\centering
\begin{tikzpicture}[scale=0.9]
    % Axes.
    \draw[->,EPICMuted] (-0.15,0) -- (3.5,0) node[right,font=\tiny]{$x$};
    \draw[->,EPICMuted] (0,-0.15) -- (0,3.1);
    % A convex f, and the chord between two points on it.
    \draw[OxfordBlue,thick,domain=0.2:2.95,smooth,variable=\x] plot ({\x},{0.32*\x*\x});
    \draw[AmazonOrange,thick] (0.6,0.1152) -- (2.9,2.6912);
    \fill[AmazonOrange] (0.6,0.1152) circle (1.5pt);
    \fill[AmazonOrange] (2.9,2.6912) circle (1.5pt);
    % The gap at x = E[X] is exactly Jensen's inequality.
    \draw[densityblue,dashed,thick] (1.75,0.98) -- (1.75,1.4032);
    \fill[AmazonOrange] (1.75,1.4032) circle (1.5pt);
    \fill[densityblue] (1.75,0.98) circle (1.5pt);
    \draw[EPICMuted,dotted] (1.75,0) -- (1.75,0.98);
    \node[font=\tiny,EPICMuted,below] at (1.75,0) {$\E[X]$};
    \node[font=\tiny,AmazonOrange,above left=-1pt and -2pt] at (1.75,1.4032) {$\E[f(X)]$};
    \node[font=\tiny,densityblue,below right=-2pt and 0pt] at (1.75,0.98) {$f(\E[X])$};
\end{tikzpicture}
\end{column}
\end{columns}
\end{frame}

% ------------------------------------------------------------------
\begin{frame}{What we have set up}
\begin{itemize}
    \item Information measures \textit{how unexpected} something was, never what it meant.
    \pause
    \item It is always measured against an assumed distribution, so it describes a \blue{model of the world}
    as much as the world.
    \pause
    \item Shannon's two problems, compression and transmission, become \blue{modelling and inference} once
    we no longer control the encoder.
    \pause
    \item The payoff is limits that hold for every method, not just the one we happen to have written.
\end{itemize}
\vspace{0.4cm}
\pause
\begin{center}
\large What remains is to turn ``how unexpected'' into \orange{a number}.
\end{center}
\end{frame}

\section{Entropy}

% ------------------------------------------------------------------
\begin{frame}{Where entropy comes from}
Write $s(p)$ for how much we learn from an event of probability $p$, and ask for three things:
\begin{itemize}
    \item $s(1) = 0$: an event we were already sure of teaches us nothing;
    \item $s$ gets smaller as $p$ grows: less likely events teach us more;
    \item $s(p_1 p_2) = s(p_1) + s(p_2)$: two independent events teach us the sum of what each one does.
\end{itemize}
\pause
\vspace{0.2cm}
Only one function does all three, whatever base of logarithm we choose:
\[
    s(p) = -\log p ,
\]
\pause
and \orange{entropy} is simply its average over the distribution,
\[
    H(X) = -\sum_x p(x) \log p(x) .
\]
\vspace{0.1cm}
\pause
\muted{So the logarithm is not a convention. It is forced on us by asking that information from independent
events add up.}
\end{frame}

% ------------------------------------------------------------------
\begin{frame}{Entropy: the definition}
\begin{block}{Definition 2.1}
The \orange{entropy} of a random variable $X$ with distribution $p$ is its average surprise,
\[
    H(X) \;=\; -\sum_x p(x)\log p(x) \;=\; \E_p\big[-\log p(X)\big].
\]
\end{block}
\pause
\begin{itemize}
    \item It is a property of the \blue{distribution}, not of any outcome we happen to observe.
    \pause
    \item $H(X) \ge 0$ always, since $p(x) \le 1$ makes every term $-\log p(x)$ non-negative.
    \pause
    \item \orange{Units.} Base $2$ gives bits, base $e$ gives nats. Duchi uses base $e$ throughout, so his
    formulas carry no stray constants; we switch to bits when we want a number to interpret.
    \pause
    \item \orange{The uniform distribution has the highest entropy.} If $X$ is uniform on $m$ outcomes then
    every outcome has surprise $\log m$, so $H(X) = \log m$; and no distribution on those outcomes does
    better, $H(X) \le \log m$, with equality only in the uniform case. (We will prove this later)
\end{itemize}
\end{frame}

% ------------------------------------------------------------------
\begin{frame}{The simplest example: a biased coin}
\begin{columns}[T,totalwidth=\textwidth]
\begin{column}{0.46\textwidth}
For $X \sim \text{Bernoulli}(p)$ we write
\[
    h_2(p) = -p\log p - (1-p)\log(1-p).
\]
\pause
\begin{itemize}
    \item Zero at $p=0$ and $p=1$: a coin that always lands the same way holds no surprise.
    \pause
    \item Largest at $p = 1/2$, where it equals exactly one bit.
    \pause
    \item \blue{Flat on top, steep at the edges.} Moving from $p=0.5$ to $p=0.6$ costs almost nothing
    ($0.971$ bits), while moving from $0.01$ to $0.001$ changes a great deal.
\end{itemize}
\end{column}
\begin{column}{0.5\textwidth}
\centering
\vspace{0.2cm}
\begin{tikzpicture}[x=3.5cm,y=1.85cm]
    \draw[->,EPICMuted] (-0.04,0) -- (1.14,0) node[right,font=\tiny]{$p$};
    \draw[->,EPICMuted] (0,-0.06) -- (0,1.25);
    \node[font=\tiny,EPICMuted,above left=-3pt] at (0,1.25) {bits};
    \draw[EPICMuted,dotted] (0.5,0) -- (0.5,1);
    \draw[OxfordBlue,thick,domain=0.004:0.996,samples=90,smooth,variable=\x]
        plot ({\x},{-\x*ln(\x)/ln(2) - (1-\x)*ln(1-\x)/ln(2)});
    \fill[AmazonOrange] (0.5,1) circle (1.5pt);
    \node[font=\tiny,AmazonOrange,above] at (0.5,1) {$1$ bit};
    \node[font=\tiny,EPICMuted,below] at (0.5,0) {$1/2$};
    \node[font=\tiny,EPICMuted,below] at (1,0) {$1$};
\end{tikzpicture}
\end{column}
\end{columns}
\end{frame}

% ------------------------------------------------------------------
\begin{frame}{Uncertainty left over: conditional entropy}
\begin{block}{Definition 2.2}
For a single observed value, $H(X \mid Y = y) = -\sum_x p(x \mid y)\log p(x \mid y)$, and averaging over $Y$,
\[
    H(X \mid Y) \;=\; \sum_y p(y)\, H(X \mid Y = y).
\]
\end{block}
\pause
\vspace{0.1cm}
\begin{itemize}
    \item Read it as: \textit{how much uncertainty about $X$ is left once we have seen $Y$}, averaged over the
    values $Y$ might take.
    \pause
    \item The averaging is not cosmetic. $H(X \mid Y)$ is one number describing the observation
    \blue{in general}, not a verdict on any single measurement.
    \pause
    \item In supervised learning $H(Y \mid X)$ is the uncertainty still in the label once the input is known:
    the part of the problem \orange{no model can remove}.
\end{itemize}
\end{frame}

% ------------------------------------------------------------------
\begin{frame}{The chain rule}
\begin{block}{Proposition 2.3 (chain rule)}
\[
    H(X_1, \dots, X_n) \;=\; H(X_1) + H(X_2 \mid X_1) + \cdots + H(X_n \mid X_1, \dots, X_{n-1}).
\]
\end{block}
\pause
\vspace{0.1cm}
\begin{itemize}
    \item In words: the uncertainty in a whole sequence is the uncertainty in the first symbol, plus whatever
    is left in the second once the first is known, and so on down the sequence.
    \pause
    \item It follows from writing $p(x,y) = p(x)\,p(y \mid x)$ and taking logarithms, which turns the product
    into a sum.
    \pause
    \item \orange{This is the factorisation autoregressive models are trained on.} The per-token cross-entropy
    loss of a language model is the sample version of exactly this sum, one term per position.
\end{itemize}
\vspace{0.3cm}
\pause
\muted{Chain rules are the workhorse of the subject. Almost every proof in Chapter 2 is an expansion of some
quantity in two different orders, followed by comparing the results.}
\end{frame}

% ------------------------------------------------------------------
\begin{frame}{Proof of Proposition 2.3}
It is enough to prove the two-variable case, then induct. Start from $p(x,y) = p(x)\,p(y \mid x)$:
\vspace{-0.1cm}
% \pause is silently dropped inside amsmath alignments, so the reveal is done with
% explicit overlay specs. Keep the & outside the \onslide group or the tab is lost.
\begin{align*}
    H(X,Y) &= -\sum_{x,y} p(x,y)\log p(x,y)
            = -\sum_{x,y} p(x)\,p(y \mid x)\,\big[\log p(x) + \log p(y \mid x)\big] \\[0.15cm]
           &\onslide<2->{= \underbrace{-\sum_x p(x)\log p(x) \sum_y p(y \mid x)}_{\textstyle \sum_y p(y \mid x) \,=\, 1}
              \;-\; \sum_x p(x) \sum_y p(y \mid x)\log p(y \mid x)} \\[0.2cm]
           &\onslide<3->{= H(X) + H(Y \mid X).}
\end{align*}
\vspace{-0.2cm}
\begin{itemize}
    \item<4-> The whole proof is: \blue{factor the joint, take logs, and the product becomes a sum}.
    \item<5-> For the general case put $X = (X_1,\dots,X_{n-1})$ and $Y = X_n$, and repeat.
\end{itemize}
\end{frame}

% ------------------------------------------------------------------
\begin{frame}{Observing can only help, on average}
\begin{block}{Proposition 2.4 (conditioning reduces entropy)}
\[
    H(X \mid Y) \;\le\; H(X).
\]
\end{block}
\pause
\vspace{0.1cm}
\begin{itemize}
    \item Looking at $Y$ cannot, on average, leave us more uncertain about $X$ than we started.
    \pause
    \item The gap between the two sides turns out to be exactly the mutual information $I(X;Y)$ of
    Section 4, so proving this also gets us that $I(X;Y) \ge 0$.
\end{itemize}
\vspace{0.25cm}
\pause
\begin{alertblock}{The average is doing the work}
A \textit{particular} observation can leave us worse off: $H(X \mid Y = y) > H(X)$ is perfectly possible.
A medical test that comes back inconclusive is the everyday version. The theorem promises only that such
outcomes are outweighed by the informative ones.
\end{alertblock}
\end{frame}

% ------------------------------------------------------------------
\begin{frame}{Proof of Proposition 2.4}
Write both entropies as sums over the joint distribution, using $p(x) = \sum_y p(x,y)$:
\vspace{-0.1cm}
\[
    H(X) - H(X \mid Y)
      \;=\; \sum_{x,y} p(x,y) \log \frac{p(x \mid y)}{p(x)}
      \;=\; \sum_{x,y} p(x,y) \log \frac{p(x,y)}{p(x)\,p(y)} .
\]
\pause
\vspace{-0.15cm}
Now negate, and apply Jensen's inequality to the concave function $\log$:
\vspace{-0.1cm}
\[
    -\big[H(X) - H(X \mid Y)\big]
      \;=\; \sum_{x,y} p(x,y) \log \frac{p(x)\,p(y)}{p(x,y)}
      \;\;\overset{\text{Jensen}}{\le}\;\;
      \log \sum_{x,y} p(x,y)\,\frac{p(x)\,p(y)}{p(x,y)}
      \;=\; \log 1 \;=\; 0 .
\]
\pause
\vspace{-0.1cm}
\begin{itemize}
    \item So $H(X) - H(X \mid Y) \ge 0$, which is the claim.
    \pause
    \item \muted{The sum collapsed because $\sum_{x,y} p(x)\,p(y) = 1$: the weights $p(x,y)$ cancel against
    the ratio, leaving the product distribution, which is still a distribution.}
\end{itemize}
\vspace{0.15cm}
\pause
\muted{This is the first place Jensen earns its keep, and the argument recurs verbatim in Section 3 to show
that the KL divergence is never negative.}
\end{frame}

% ------------------------------------------------------------------
\begin{frame}{The whole is no more uncertain than its parts}
\begin{block}{Proposition 2.5 (subadditivity)}
\[
    H(X_1, \dots, X_n) \;\le\; H(X_1) + H(X_2) + \cdots + H(X_n),
\]
with equality exactly when $X_1, \dots, X_n$ are mutually independent.
\end{block}
\pause
Chain rule, then drop each conditioning in turn:
\vspace{-0.15cm}
\begin{align*}
    H(X_1, \dots, X_n) &\;\overset{\text{Prop. }2.3}{=}\; \sum_{i=1}^n H(X_i \mid X_1, \dots, X_{i-1}) \\[0.05cm]
                       &\;\overset{\text{Prop. }2.4}{\le}\; \sum_{i=1}^n H(X_i) .
\end{align*}
\pause
\vspace{-0.28cm}
\begin{itemize}
    \item Equality needs every step tight, so each $X_i$ is independent of all those before it: mutual
    independence.
    \pause
    \item \orange{Describing variables together never costs more than one by one.} The saving is their
    dependence: at $n = 2$ the gap $H(X) + H(Y) - H(X,Y)$ is the mutual information of Section 4.
\end{itemize}
\end{frame}

% ------------------------------------------------------------------
\begin{frame}{Entropy in machine learning}
\begin{itemize}
    \item \orange{Predictive uncertainty.} The entropy of a softmax output is the standard one-number summary
    of how unsure a classifier is.
    \pause
    \item \orange{Perplexity} is $e^{H}$, the effective number of equally likely choices. A language model at
    $H = 2.08$ nats per token is as uncertain as if it faced $8$ equally likely tokens at every step.
    \pause
    \item \orange{The error floor.} $H(Y \mid X)$ is the part of the loss that no model can remove. A
    cross-entropy of $2.1$ nats is not evidence of a bad model until we know what the floor is.
\end{itemize}
\vspace{0.25cm}
\pause
\begin{alertblock}{Entropy cannot see the outcomes}
$H$ depends only on the probabilities, so relabelling or reordering the classes leaves it unchanged. It is
therefore a poor measure of uncertainty when the outcomes are ordered or numerical, where being wrong by a
little is not the same as being wrong by a lot.
\end{alertblock}
\end{frame}

% ------------------------------------------------------------------
\begin{frame}{Exercise 2.1}
\begin{exampleblock}{Work in bits}
A weather model outputs $p = (0.7,\, 0.2,\, 0.1)$ over (sun, cloud, rain).
\vspace{0.15cm}
\begin{enumerate}
    \item Compute $H(p)$.
    \item Compute it for $u = (\tfrac13,\tfrac13,\tfrac13)$ and for $c = (0.98,\, 0.01,\, 0.01)$,
    and put the three in order.
    \item The model outputs $c$ on an input from a class it has never seen before, so its entropy is tiny.
    What is entropy failing to measure?
\end{enumerate}
\end{exampleblock}
\vspace{0.3cm}
\pause
\muted{Useful numbers: $\log_2 0.7 = -0.515$, $\log_2 0.2 = -2.322$, $\log_2 0.1 = -3.322$,
$\log_2 3 = 1.585$, $\log_2 0.98 = -0.029$.}
\end{frame}

% ------------------------------------------------------------------
\begin{frame}{Exercise 2.1, solution}
\begin{itemize}
    \item[(1)] $H(p) = 0.360 + 0.464 + 0.332 = \orange{1.157}$ bits.
    \pause
    \item[(2)] $H(u) = \log_2 3 = \orange{1.585}$ bits; $H(c) = 0.029 + 0.133 = \orange{0.161}$ bits.
    So $H(c) < H(p) < H(u)$: the confident model has least entropy, the uniform one most.
    \pause
    \item[(3)] Entropy measures only the \blue{spread of the prediction}, which mixes two different
    situations: the world being genuinely ambiguous, and the model extrapolating confidently into territory
    it has never seen.
\end{itemize}
\vspace{0.25cm}
\pause
\begin{alertblock}{Where this is going}
Separating those two requires a distribution \textit{over models}, not a single prediction. The quantity that
does it is mutual information, in Section 4.
\end{alertblock}
\end{frame}

% ------------------------------------------------------------------
\section{KL Divergence}

% ------------------------------------------------------------------
\begin{frame}{Why one distribution is not enough}
Entropy scores \textit{one} distribution. Modelling always involves \blue{two}: the truth $P$, and whatever
model $Q$ we have fitted.
\pause
\vspace{0.25cm}
\begin{itemize}
    \item Suppose we build our code believing the data follow $Q$, so a symbol $x$ costs $-\log q(x)$ to send.
    \pause
    \item The symbols actually arrive from $P$, so our average cost is $\E_P[-\log q(X)]$.
    \pause
    \item The best we could have done, knowing $P$, was $H(P)$.
\end{itemize}
\vspace{0.3cm}
\pause
\begin{center}
\large The \orange{excess} we pay for believing the wrong model is what we now measure.
\end{center}
\end{frame}

% ------------------------------------------------------------------
\begin{frame}{The Kullback-Leibler divergence}
\begin{block}{Definition 3.1}
For distributions $P$ and $Q$ on the same set, $\KL{P}{Q} = \sum_x p(x)\log\frac{p(x)}{q(x)}$.
Since $\log\frac{p}{q} = \log p - \log q$, this splits as
\vspace{-0.15cm}
\[
    \KL{P}{Q} \;=\; \underbrace{\left(-\sum_x p(x) \log q(x)\right)}_{\textstyle \text{cross-entropy}}
      \;-\; \underbrace{\left(-\sum_x p(x) \log p(x)\right)}_{\textstyle H(P)} .
\]
\end{block}
\pause
\vspace{0.05cm}
\begin{itemize}
    \item Equivalently $\E_P\!\left[\log \frac{p(X)}{q(X)}\right]$. The expectation is under $P$ alone, a
    first hint that the two arguments are not symmetric.
    \pause
    \item \orange{So the KL divergence is cross-entropy minus entropy.} Only the first bracket involves $Q$,
    so it is the only part a model can change.
    \pause
    \item Conventions: $0\log 0 = 0$, and $\KL{P}{Q} = +\infty$ once $q(x) = 0$ where $p(x) > 0$: ruling out
    something that then happens is infinitely bad.
\end{itemize}
\end{frame}

% ------------------------------------------------------------------
\begin{frame}{The KL divergence is never negative}
\begin{block}{Proposition 3.2}
$\KL{P}{Q} \ge 0$ for all $P, Q$, with equality exactly when $P = Q$.
\end{block}
\pause
\vspace{0.1cm}
Negate, and apply Jensen's inequality to the concave function $\log$:
\vspace{-0.1cm}
\[
    -\KL{P}{Q} \;=\; \E_P\!\left[\log \frac{q(X)}{p(X)}\right]
    \;\;\overset{\text{Jensen}}{\le}\;\;
    \log\, \E_P\!\left[\frac{q(X)}{p(X)}\right]
    \;=\; \log \sum_x p(x)\,\frac{q(x)}{p(x)}
    \;=\; \log 1 \;=\; 0 .
\]
\pause
\vspace{-0.1cm}
\begin{itemize}
    \item Since $\log$ is \textit{strictly} concave, equality needs $q(X)/p(X)$ constant, which forces $P = Q$.
    \pause
    \item \muted{This is the same computation as the proof of Proposition 2.4. That proposition was this one
    in disguise, applied to the joint distribution against the product of its marginals.}
\end{itemize}
\end{frame}

% ------------------------------------------------------------------
\begin{frame}{Uniform has Maximum Entropy}
\begin{block}{Proposition 3.3 (uniform maximises entropy)}
If $X$ takes values in a set of size $m$, then $H(X) \le \log m$, with equality exactly when $X$ is uniform.
\end{block}
\pause
\vspace{0.1cm}
Let $U$ be uniform, so $q(x) = 1/m$ for every $x$. Then
\vspace{-0.1cm}
\begin{align*}
    0 \;\overset{\text{Prop. }3.2}{\le}\; \KL{P}{U}
      &\;=\; \sum_x p(x)\log\frac{p(x)}{1/m}
       \;=\; \sum_x p(x)\log p(x) \;+\; \log m \sum_x p(x) \\
      &\;=\; -H(X) + \log m .
\end{align*}
\pause
\vspace{-0.1cm}
\begin{itemize}
    \item Rearranging gives $H(X) \le \log m$, and equality holds exactly when $P = U$.
\end{itemize}
\end{frame}

% ------------------------------------------------------------------
\begin{frame}{What training actually minimises}
\begin{block}{Proposition 3.4 (cross-entropy decomposition)}
\[
    \underbrace{\E_P[-\log q(X)]}_{\text{cross-entropy loss}}
    \;=\; \underbrace{H(P)}_{\text{not our fault}}
    \;+\; \underbrace{\KL{P}{Q}}_{\text{our fault}} .
\]
\end{block}
\pause
\vspace{0.05cm}
The proof is to add and subtract $\log p(x)$:
$\E_P[-\log q] = \E_P[-\log p] + \E_P[\log p/q]$.
\pause
\vspace{0.15cm}
\begin{itemize}
    \item \orange{Training with cross-entropy is minimising a KL divergence} to the data distribution. Only the
    second term depends on the model.
    \pause
    \item The loss can never fall below $H(P)$. A cross-entropy of $2.1$ nats says nothing on its own.
    \pause
    \item Taking $P$ to be the \blue{empirical} distribution of the sample makes this exactly
    \orange{maximum likelihood}. The two are the same optimisation written differently.
\end{itemize}
\end{frame}

% ------------------------------------------------------------------
\begin{frame}{A divergence, not a distance}
\begin{itemize}
    \item $\KL{P}{Q} \ne \KL{Q}{P}$ in general, and there is no triangle inequality.
    \pause
    \item So ``KL distance'' is a misnomer. The right word is \orange{divergence}: it vanishes exactly when
    the two distributions agree, but it does not make the space of distributions a metric space.
    \pause
    \item Vanishing exactly at $P = Q$ is all we need for it to work as a \textit{training objective}, which
    is why the asymmetry so rarely gets mentioned.
    \pause
    \item \muted{Part 2 introduces quantities that genuinely are metrics, among them the total variation and
    Hellinger distances, and explains when the difference starts to matter.}
\end{itemize}
\vspace{0.3cm}
\pause
\begin{center}
\large But the asymmetry is not a technicality. \orange{The two directions fit different models.}
\end{center}
\end{frame}

% ------------------------------------------------------------------
\begin{frame}{Forward and reverse KL fit different models}
\begin{columns}[T,totalwidth=\textwidth]
\begin{column}{0.5\textwidth}
\vspace{0.1cm}
\begin{itemize}
    \item \orange{Forward}, $\KL{P}{Q}$: heavily penalised where $p$ is large and $q$ small, so $Q$ must
    \blue{cover} every mode. It spreads out, putting mass in the valley where $P$ has none.
    \pause
    \item \orange{Reverse}, $\KL{Q}{P}$: heavily penalised where $q$ is large and $p$ small, so $Q$
    \blue{retreats} into one mode and stays narrow.
    \pause
    \item Maximum likelihood minimises the forward direction; variational inference minimises the reverse.
\end{itemize}
\end{column}
\begin{column}{0.46\textwidth}
\centering
\vspace{0.3cm}
\begin{tikzpicture}[x=0.80cm,y=1.35cm]
    \draw[->,EPICMuted] (-3.5,0) -- (3.6,0) node[right,font=\tiny]{$x$};
    % The target P: two well separated modes.
    \draw[OxfordBlue,thick,domain=-3.4:3.4,samples=140,smooth,variable=\x]
        plot ({\x},{0.5*exp(-(\x+1.5)^2/0.5) + 0.5*exp(-(\x-1.5)^2/0.5)});
    % Forward KL: a single wide Gaussian covering both modes.
    \draw[AmazonOrange,thick,dashed,domain=-3.4:3.4,samples=140,smooth,variable=\x]
        plot ({\x},{0.40*exp(-(\x)^2/5.0)});
    % Reverse KL: a narrow Gaussian sitting on one mode.
    \draw[deepgreen,thick,dash dot,domain=-3.4:3.4,samples=140,smooth,variable=\x]
        plot ({\x},{0.52*exp(-(\x-1.5)^2/0.29)});
    \node[font=\tiny,OxfordBlue,above] at (-1.5,0.5) {$P$};
    \node[font=\tiny,AmazonOrange,above] at (0,0.4) {forward};
    \node[font=\tiny,deepgreen,above right=-1pt and -3pt] at (2.3,0.18) {reverse};
\end{tikzpicture}
\end{column}
\end{columns}
\vspace{0.02cm}
\pause
\begin{alertblock}{Why this matters for uncertainty}
Variational inference minimises the \textit{reverse} direction, so its posterior hides inside one mode and
reports too little uncertainty. That is a property of the objective, not a failure of the optimiser.
\end{alertblock}
\end{frame}

% ------------------------------------------------------------------
\begin{frame}{Continuous variables, and a warning}
For densities the sum becomes an integral, $\KL{P}{Q} = \int p(x)\log\frac{p(x)}{q(x)}\,dx$, and the
\blue{differential entropy} is $h(X) = -\int p(x)\log p(x)\,dx$.
\pause
\vspace{0.2cm}
\begin{alertblock}{Differential entropy is not entropy}
It can be \textit{negative}; it is not the limit of the discrete entropy; and it \textit{changes} if we
rescale or reparameterise $X$. It is not a count of bits.
\end{alertblock}
\pause
\vspace{0.15cm}
\begin{itemize}
    \item \orange{The KL divergence and the mutual information do not suffer from this}: both are unchanged
    by any invertible transformation of the variable. Report those, and treat a quoted differential entropy
    with suspicion.
    \pause
    \item For two Gaussians with the same covariance, $P = \mathsf{N}(\mu_1,\Sigma)$ and
    $Q = \mathsf{N}(\mu_2,\Sigma)$,
    \[
        \KL{P}{Q} \;=\; \tfrac12 (\mu_1-\mu_2)^\top \Sigma^{-1} (\mu_1-\mu_2),
    \]
    a squared Mahalanobis distance. \muted{This is why so many statistical rates come out proportional to
    $\Delta^2/\sigma^2$, as Part 2 will show.}
\end{itemize}
\end{frame}

% ------------------------------------------------------------------
\begin{frame}{Exercise 3.1}
\begin{enumerate}
    \item Prove Proposition 3.4, that $\E_p[-\log q] = H(p) + \KL{p}{q}$.
    \pause
    \item Let $p$ be the empirical distribution of a sample $x_1,\dots,x_n$, so $p$ puts mass $1/n$ on each
    point. Write out $\E_p[-\log q_\theta]$ and show that minimising it over $\theta$ is maximum likelihood.
    \pause
    \item Let $P = \tfrac12 \mathsf{N}(-3,1) + \tfrac12 \mathsf{N}(3,1)$ and $Q_\sigma = \mathsf{N}(0,\sigma^2)$.
    Without computing anything, say what $\arg\min_\sigma \KL{P}{Q_\sigma}$ looks like, and how
    $\arg\min_\sigma \KL{Q_\sigma}{P}$ differs. Sketch both.
\end{enumerate}
\end{frame}

% ------------------------------------------------------------------
\begin{frame}{Exercise 3.1, solution}
\begin{itemize}
    \item[(1)] $\E_p[-\log q] = \E_p[-\log p] + \E_p\!\left[\log \frac{p}{q}\right] = H(p) + \KL{p}{q}$.
    \pause
    \item[(2)] It reads $-\frac{1}{n}\sum_{i=1}^n \log q_\theta(x_i)$, the negative average log-likelihood.
    The term $H(p)$ does not involve $\theta$, so the two problems have the same solution.
    \pause
    \item[(3)] \orange{Forward} blows up wherever $p > 0$ and $q \approx 0$, so $Q$ must stretch to cover both
    modes: a wide Gaussian straddling the empty valley between them, a poor fit everywhere but a finite one.
    \orange{Reverse} blows up wherever $q > 0$ and $p \approx 0$, so $Q$ hides inside a single mode: a narrow
    Gaussian at $+3$, or equally at $-3$.
\end{itemize}
\vspace{0.2cm}
\pause
\begin{alertblock}{The point of part (3)}
Choosing a direction of KL is a modelling decision about whether you would rather be \textit{over-confident}
or \textit{over-cautious}.
\end{alertblock}
\end{frame}

% ------------------------------------------------------------------
\section{Mutual Information}

% ------------------------------------------------------------------
\begin{frame}{How much does an observation tell us?}
So far: Section 2 scored \textit{one} distribution, Section 3 compared \textit{two}.
\pause
\vspace{0.25cm}
\begin{itemize}
    \item Before looking at $Y$, our uncertainty about $X$ is $H(X)$.
    \pause
    \item After looking at $Y$, what is left is $H(X \mid Y)$.
    \pause
    \item The difference is what the observation was worth.
\end{itemize}
\vspace{0.3cm}
\pause
\begin{center}
\large That difference has a name, and it turns out to be \orange{a KL divergence in disguise}.
\end{center}
\end{frame}

% ------------------------------------------------------------------
\begin{frame}{Mutual information}
\begin{block}{Definition 4.1}
The \orange{mutual information} between $X$ and $Y$ is the divergence between their joint distribution and
the product of their marginals,
\vspace{-0.1cm}
\[
    I(X;Y) \;=\; \KL{P_{XY}}{P_X \! \times \! P_Y}
           \;=\; \sum_{x,y} p(x,y) \log \frac{p(x,y)}{p(x)\,p(y)} .
\]
\end{block}
\pause
\vspace{0.05cm}
\begin{itemize}
    \item The product $P_X \! \times \! P_Y$ is what the joint \textit{would} be if $X$ and $Y$ were
    independent. So mutual information measures \blue{how far the pair is from independence}.
    \pause
    \item Because it is a KL divergence, Proposition 3.2 applies at once: $I(X;Y) \ge 0$, with equality
    exactly when $X$ and $Y$ are independent.
\end{itemize}
\end{frame}

% ------------------------------------------------------------------
\begin{frame}{Three ways to read the same quantity}
\begin{block}{Proposition 4.2}
\[
    I(X;Y) \;=\; H(X) - H(X \mid Y) \;=\; H(Y) - H(Y \mid X).
\]
\end{block}
\pause
\vspace{0.05cm}
Using $\dfrac{p(x,y)}{p(x)\,p(y)} = \dfrac{p(x \mid y)}{p(x)}$ and splitting the logarithm,
\vspace{-0.1cm}
\[
    I(X;Y) \;=\; \underbrace{\sum_{x,y} p(x,y)\log p(x \mid y)}_{\textstyle -H(X \mid Y)}
             \;-\; \underbrace{\sum_{x,y} p(x,y)\log p(x)}_{\textstyle -H(X)} ,
\]
and the definition is symmetric in $x$ and $y$, which gives the second form.
\pause
\vspace{0.02cm}
\begin{itemize}
    \item \orange{What $Y$ says about $X$ equals what $X$ says about $Y$}, which is far from obvious in words.
    \pause
    \item \orange{We already know $I(X;Y) \ge 0$}, from Definition 4.1 and Proposition 3.2. Put that together
    with the identity above and it reads $H(X \mid Y) \le H(X)$: we have recovered \blue{Proposition 2.4}.
    \muted{Its direct proof in Section 2 was this same computation, run before we had the name for it.}
\end{itemize}
\end{frame}

% ------------------------------------------------------------------
\begin{frame}{Once we already know something}
\begin{block}{Definition 4.3 (conditional mutual information)}
Compute the mutual information separately inside each slice $Z = z$, then average:
\vspace{-0.1cm}
\[
    I(X;Y \mid Z) \;=\; \sum_z p(z)\, I(X;Y \mid Z = z)
                  \;=\; H(X \mid Z) - H(X \mid Y, Z).
\]
\end{block}
\pause
\vspace{0.05cm}
\begin{itemize}
    \item Read it as: \blue{how much $Y$ still tells us about $X$ once $Z$ is already known}. Every reading
    from Proposition 4.2 carries over, with $Z$ held fixed throughout.
    \pause
    \item Each slice is a mutual information, hence non-negative, so $I(X;Y \mid Z) \ge 0$ as well.
\end{itemize}
\vspace{0.0cm}
\pause
\begin{alertblock}{Extra context can raise information, unlike entropy}
$I(X;Y \mid Z)$ may \textit{exceed} $I(X;Y)$. Let $X$ and $Y$ be independent fair bits and $Z = X \oplus Y$.
Then $I(X;Y) = 0$, yet knowing $Z$ makes each bit determine the other, so $I(X;Y \mid Z) = 1$ bit. Contrast
Proposition 2.4: conditioning always reduces \textit{entropy} on average, but it does not always reduce
\textit{information}.
\end{alertblock}
\end{frame}

% ------------------------------------------------------------------
\begin{frame}{Information adds up, but not twice over}
\begin{block}{Proposition 4.4 (chain rule for information)}
\[
    I(X; Y_1,\dots,Y_n) \;=\; \sum_{i=1}^n I(X; Y_i \mid Y_1, \dots, Y_{i-1}).
\]
\end{block}
\pause
\vspace{0.1cm}
\begin{itemize}
    \item Each observation contributes what it adds \textit{given everything already seen}. The proof is the
    entropy chain rule of Proposition 2.3, applied twice and subtracted.
    \pause
    \item If the $Y_i$ are independent given $X$, conditioning can only shrink each term, so
    \[
        I(X; Y_1,\dots,Y_n) \;\le\; \sum_{i=1}^n I(X; Y_i).
    \]
    \pause
    \item \orange{$n$ observations are worth at most $n$ times one, and usually less.} The shortfall is
    \blue{redundancy}: the second measurement repeats part of what the first already said.
\end{itemize}
\vspace{0.1cm}
\pause
\muted{Hence batch-aware acquisition scores the joint information, rather than summing individual scores.}
\end{frame}

% ------------------------------------------------------------------
\begin{frame}{Processing cannot create information}
\begin{block}{Proposition 4.5 (data processing inequality)}
If $X \to Y \to Z$ is a Markov chain, meaning $Z$ depends on $X$ only through $Y$, then
\[
    I(X;Z) \;\le\; I(X;Y).
\]
\end{block}
\pause
\vspace{0.1cm}
Expand $I(X; Y,Z)$ in the two possible orders using Proposition 4.4:
\vspace{-0.1cm}
\begin{align*}
    I(X; Y,Z) &\;=\; I(X;Z) + I(X;Y \mid Z) \\[0.15cm]
              &\;=\; I(X;Y) + \underbrace{I(X;Z \mid Y)}_{\textstyle =\,0} .
\end{align*}
\pause
\vspace{-0.1cm}
The marked term vanishes because $X$ and $Z$ are independent given $Y$. Since $I(X;Y \mid Z) \ge 0$, the
first line is at least $I(X;Z)$, which gives the claim.
\end{frame}

% ------------------------------------------------------------------
\begin{frame}{What that rules out}
\begin{itemize}
    \item \orange{No architecture creates information.} If the input does not determine the label, no network
    on top of it can. Depth makes information \blue{accessible}, never more abundant.
    \pause
    \item \orange{Post-processing is free of charge and free of benefit.} Temperature scaling, Platt scaling
    and taking an argmax are all functions of the logits, so none of them can raise $I(Y; \cdot)$. They improve
    how probabilities are \textit{reported}, not what is known.
    \pause
    \item \orange{A diagnostic.} If a frozen backbone yields features $Z$ with $I(Y;Z)$ small, then no head,
    however large, will rescue the accuracy. The thing to change is the representation or the data.
    \pause
    \item \muted{The information bottleneck turns this into an objective: deliberately shrink $I(X;Z)$ while
    holding $I(Z;Y)$ up, to keep only what the label needs.}
\end{itemize}
\end{frame}

% ------------------------------------------------------------------
\begin{frame}{Splitting uncertainty in two}
Let $\theta$ index the model (a posterior sample, or a member of an ensemble), $x$ a fixed input, $Y$ the label.
\vspace{0.1cm}
\pause
\begin{block}{The decomposition}
\[
    \underbrace{I(Y;\theta \mid x)}_{\text{epistemic}}
    \;=\; \underbrace{H\big(\E_{\theta}[\,p(y \mid x,\theta)\,]\big)}_{\text{total predictive uncertainty}}
    \;-\; \underbrace{\E_{\theta}\big[H(p(y \mid x,\theta))\big]}_{\text{aleatoric}} .
\]
\end{block}
\pause
\vspace{0.1cm}
\begin{itemize}
    \item This is Proposition 4.2 applied to $Y$ and $\theta$, and nothing more.
    \pause
    \item \orange{Read it as a question about learning.} Epistemic uncertainty is how much a label at $x$ would
    teach us about $\theta$. If seeing it would change our beliefs about the model, we are unsure of the model.
    \pause
    \item Whatever uncertainty survives after averaging over every model is \blue{noise in the world}, and no
    amount of data removes it.
\end{itemize}
\end{frame}

% ------------------------------------------------------------------
\begin{frame}{Exercise 4.1}
\begin{exampleblock}{Two models, one input. Work in bits}
Fix $x$. The posterior over models is two equally likely parameters $\theta_1, \theta_2$, and the label is
binary. Write $p_i = p(y{=}1 \mid x, \theta_i)$ for what model $i$ predicts.
\vspace{0.1cm}
\begin{center}
\footnotesize
\begin{tabular}{ccc}
\toprule
\textbf{Case} & $p_1$ & $p_2$ \\
\midrule
A & $0.5$ & $0.5$ \\
B & $1.0$ & $0.0$ \\
C & $0.9$ & $0.1$ \\
\bottomrule
\end{tabular}
\end{center}
\vspace{0.05cm}
\begin{enumerate}
    \item Fill in the four columns below for each case.
    \item Which cases can a single averaged softmax not tell apart?
    \item One label can be bought: spend it on an input of type A, or of type B?
\end{enumerate}
\end{exampleblock}
\end{frame}

% ------------------------------------------------------------------
\begin{frame}{Exercise 4.1, how to compute it}
Everything is built from the binary entropy function
$h_2(p) = -p\log_2 p - (1-p)\log_2(1-p)$, in three steps.
\pause
\vspace{0.1cm}
\begin{block}{The recipe}
\vspace{-0.35cm}
\begin{align*}
    \text{1. average the two predictions:} \quad
        \bar p &\;=\; \tfrac12 p_1 + \tfrac12 p_2 \\[0.1cm]
    \text{2. entropy of that average:} \quad
        H(\bar p) &\;=\; h_2(\bar p)
        \qquad\qquad \text{\muted{(total uncertainty)}} \\[0.1cm]
    \text{3. average of the two entropies:} \quad
        \E_\theta[H] &\;=\; \tfrac12 h_2(p_1) + \tfrac12 h_2(p_2)
        \quad\; \text{\muted{(aleatoric)}}
\end{align*}
\vspace{-0.35cm}
\[
    \text{and then} \qquad I(Y;\theta \mid x) \;=\; H(\bar p) - \E_\theta[H] .
\]
\end{block}
\pause
\vspace{0.1cm}
\muted{The only values you need: $h_2(0.5) = 1$, $\; h_2(0) = h_2(1) = 0$, $\;$ and
$h_2(0.9) = h_2(0.1) = 0.469$.}
\end{frame}

% ------------------------------------------------------------------
\begin{frame}{Exercise 4.1, case C worked through}
Take $p_1 = 0.9$ and $p_2 = 0.1$. First the binary entropy of a single model's prediction:
\vspace{-0.1cm}
\[
    h_2(0.9) \;=\; -0.9\log_2 0.9 \;-\; 0.1\log_2 0.1 \;=\; 0.137 + 0.332 \;=\; 0.469 ,
\]
and $h_2(0.1)$ is the same by symmetry. Now the three steps:
\pause
\vspace{-0.1cm}
\begin{align*}
    \bar p                 &\;=\; \tfrac12 (0.9) + \tfrac12 (0.1) \;=\; 0.5 \\[0.1cm]
    H(\bar p)              &\;=\; h_2(0.5) \;=\; 1.000 \text{ bits}
                              &&\text{\muted{total}} \\[0.1cm]
    \E_\theta[H]           &\;=\; \tfrac12 (0.469) + \tfrac12 (0.469) \;=\; 0.469 \text{ bits}
                              &&\text{\muted{aleatoric}} \\[0.1cm]
    I(Y;\theta \mid x)     &\;=\; 1.000 - 0.469 \;=\; \mathbf{0.531} \text{ bits}
                              &&\text{\muted{epistemic}}
\end{align*}
\pause
\vspace{-0.1cm}
\muted{Cases A and B go the same way. In A both models say $0.5$, so step 3 gives $1$ and the difference is
$0$; in B the models say $1$ and $0$, so each has $h_2 = 0$, step 3 gives $0$, and the difference is a full bit.}
\end{frame}

% ------------------------------------------------------------------
\begin{frame}{Exercise 4.1, solution}
\begin{center}
\footnotesize
\begin{tabular}{ccccc}
\toprule
\textbf{Case} & averaged $\bar p$ & $H(\bar p)$ & $\E_\theta[H]$ & $I(Y;\theta \mid x)$ \\
\midrule
A & $0.5$ & $1.000$ & $1.000$ & $\mathbf{0.000}$ \\
B & $0.5$ & $1.000$ & $0.000$ & $\mathbf{1.000}$ \\
C & $0.5$ & $1.000$ & $h_2(0.9) = 0.469$ & $\mathbf{0.531}$ \\
\bottomrule
\end{tabular}
\end{center}
\pause
\vspace{0.15cm}
\begin{itemize}
    \item[(2)] \orange{All three.} Every case has predictive entropy of exactly one bit. A single averaged
    softmax cannot distinguish a genuinely even coin from two confident models in flat contradiction.
    \pause
    \item[(3)] Type B. There the label is worth a full bit \textit{about $\theta$}, while in A it is a coin
    toss that teaches us nothing about the model.
\end{itemize}
\vspace{0.15cm}
\pause
\begin{alertblock}{The moral}
Predictive entropy tells us \textit{that} a model is unsure. Mutual information tells us \textit{whether that
is the world's fault or the model's}, and only the second kind can be fixed by collecting more data.
\end{alertblock}
\end{frame}

% ------------------------------------------------------------------
\begin{frame}{Exercise 4.2 \muted{(Duchi, Exercise 2.15a)}}
\begin{exampleblock}{Ensemble Setting}
Let $\theta$ take one of $K$ values $\theta_1, \dots, \theta_K$ with weights $w_1, \dots, w_K$, and write
$p_k(y)$ for what model $k$ predicts at a fixed input $x$. Let $\bar p = \sum_k w_k\, p_k$ be the ensemble
mean. Show that
\vspace{-0.1cm}
\[
    I(Y;\theta) \;=\; \sum_{k=1}^K w_k \, \KL{p_k}{\bar p} .
\]
\end{exampleblock}
\pause
\vspace{0.15cm}
\begin{itemize}
    \item Everything needed is Definition 4.1, Definition 3.1, and a little algebra.
    \pause
    \item \muted{Hint: write out the joint distribution of $(\theta, Y)$ and the product of its marginals,
    then watch what happens to the weights $w_k$.}
\end{itemize}
\end{frame}

% ------------------------------------------------------------------
\begin{frame}{Exercise 4.2, solution}
The joint is $p(\theta_k, y) = w_k\, p_k(y)$; the marginals are $w_k$ and $\bar p(y)$. By Definition 4.1,
\vspace{-0.1cm}
% \pause is silently dropped inside amsmath alignments, so the reveal uses explicit
% overlay specs. Keep the & outside the \onslide group or the tab is lost.
\begin{align*}
    I(Y;\theta) &\;=\; \sum_{k}\sum_{y} w_k\, p_k(y)
                   \log \frac{w_k\, p_k(y)}{w_k\, \bar p(y)} \\[0.12cm]
                &\onslide<2->{\;=\; \sum_{k} w_k \sum_{y} p_k(y) \log \frac{p_k(y)}{\bar p(y)}} \\[0.12cm]
                &\onslide<3->{\;=\; \sum_{k} w_k \, \KL{p_k}{\bar p} .}
\end{align*}
\vspace{-0.2cm}
\begin{itemize}
    \item<4-> The weights cancel because they appear in the joint \textit{and} in the product of marginals.
    What survives is the distance from each model to the consensus.
    \item<5-> \orange{So epistemic uncertainty is disagreement, measured in KL.} By Proposition 3.2 it is zero
    exactly when every model predicts the same thing.
    \item<6-> \muted{Check against case C: $\KL{\text{Bern}(0.9)}{\text{Bern}(0.5)} = 0.531$ bits, and both
    models sit that far from the mean, so the average is $0.531$ bits: the same number by a different route.}
\end{itemize}
\end{frame}

% ------------------------------------------------------------------
\begin{frame}{What Part 2 will do with all this}
We now have the three quantities and their basic properties. Part 2 puts them to work.
\vspace{0.2cm}
\pause
\begin{itemize}
    \item \orange{Other ways to compare distributions.} KL is one member of a family, the $f$-divergences.
    Some of that family are genuine \blue{distances}, stay bounded, and survive the support mismatch that
    sends KL to infinity.
    \pause
    \item \orange{How well can anyone tell two distributions apart?} Le Cam's inequality says the answer is
    exactly the total variation distance, which turns ``how much data do I need'' into a calculation.
    \pause
    \item \orange{How well can anyone classify?} Fano's inequality converts the mutual information carried by
    a representation into a \blue{floor on the error rate} that no method can beat.
    \pause
    \item \orange{What entropy actually means.} The source coding theorem: the best possible code for a source
    has average length between $H(X)$ and $H(X) + 1$. Entropy stops being an analogy for bits and becomes a
    count of them.
\end{itemize}
\vspace{0.25cm}
\pause
\muted{Everything there rests on what we proved today, and on Jensen's inequality, which by now has done most
of the work three times over.}
\end{frame}

% ------------------------------------------------------------------

% ------------------------------------------------------------------
% References. Uncomment together with the natbib block in the preamble.
% \begin{frame}[allowframebreaks]{References}
% \footnotesize
% \bibliographystyle{apalike}
% \bibliography{reference}
% \end{frame}

% ------------------------------------------------------------------
\begin{frame}[plain]
\begin{center}
\vspace{2cm}
{\Huge\orange{Thank you!}}\\[0.4cm]
{\Large Any Questions?}
\end{center}
\end{frame}

\end{document}
